% !Mode:: "TeX:UTF-8"
% 中英文摘要，照《研究生学位论文书写范例（理工类）》。一个 abstract 环境排一篇，写在哪里就排在哪里；
% 英文那篇写 [lang=en]。关键词写在封面文件的 info 里（keywords），用英文逗号分隔，分号由模板加；
% 也可以写在环境的可选参数里：\begin{abstract}[keywords={甲, 乙}]。
%
% 中文版与英文版共用这个文件。\abbr 按它所在的那篇摘要的语种排：中文摘要里是
% “有限元方法（Finite Element Method，FEM）”，英文摘要里是“Finite Element Method (FEM)”，
% 英文版也一样。每篇摘要排完“用过”的状态清一次，下一篇摘要与正文头一次出现都全写。
% \hitindex 写什么排什么，含汉字的进中文索引。

\begin{abstract}
\hitindex{气体静压轴承}由于具有运动精度高、摩擦损耗小、发热变形小、寿命长、无污染等特点，在航空航天工业、半导体工业、纺织工业和测量仪器中得到广泛应用。本文在分析国内外气体静压轴承的基础上，以改善气体静压轴承的静态特性和稳定性为目的，通过理论分析、仿真计算和实验研究对局部多孔质气体静压止推轴承进行了研究，同时分析轴承的结构参数和工作参数对局部多孔质气体静压止推轴承工作特性的影响，为局部多孔质气体静压轴承的设计和工程应用奠定理论基础。

建立基于分形几何理论的\hitindex{多孔质石墨}渗透率与分形维数之间关系的数学模型，该模型可预测多孔质石墨的渗透率，并可直观描述孔隙的大小对渗透率的影响。

本文在理论分析的基础上，建立局部多孔质气体静压止推轴承静态特性的数学模型，在此基础上，通过工程方法和\abbr{FEM}对所建立的模型进行求解。在采用\abbr{FEM}时，首先通过加权余量法，将二阶偏微分方程降为一阶，从而，降低了对插值函数连续度的要求，同时，方便采用有限元方法进行求解。

……
\end{abstract}

\begin{abstract}[lang=en]
Externally pressurized gas bearing has been widely used in the field of aviation, semiconductor, weave, and measurement apparatus because of its advantage of high accuracy, little friction, low heat distortion, long life-span, and no pollution. In this thesis, based on the domestic and overseas researching development about externally pressurized gas bearing, the author investigated the partial porous externally pressurized gas thrust bearing by theoretical analysis, computer simulation, and experiments to improve its static characteristics and stability. The effects of structure and operating parameters on partial porous externally pressurized gas bearing has been studied. Therefore, a theoretical foundation for the designing and application for the partial porous externally pressurized gas bearing has been presented.

Based on the fractal theory, a model was established to demonstrate the relationship between the porous graphite permeability and the fractal dimension. It can predict the permeability of porous graphite and show the effects of the pore size on the permeability.

In this thesis, the author established a model about the static characteristics of partial porous externally pressurized gas thrust bearing, and it was analyzed by engineering solution and \abbr{FEM}. While using \abbr{FEM}, the second-order partial differential equation was reduced to one-order by adopting Galerkin weighted residual method, for decreasing the continuity degree requirement of the interpolation function and facilitating to the calculation.

…
\end{abstract}
